\begin{textsample}{NEG Dim 4 – global\_north\_2024\_10 – Score -1.00 – global\_north\_2024\_10/115867607.txt}  \label{ex:f4_neg_054}
One year of war in Gaza -- protect journalists now, says IPI This week marked the grim \textbf{one-year} anniversary of the surprise October 7 Hamas attack on Israel and the beginning of the Israeli war on Gaza -- a conflict that has taken a devastating toll on journalists and media outlets in Palestine, reports the International Press Institute.
In Gaza, Israeli strikes have killed at least 123 journalists ( Gaza media sources say 178 killed ) -- the largest number of journalists to be killed in any armed conflict in this span of time to date.
Dozens of media outlets have been leveled.
Independent investigations such as those conducted by Forbidden Storieshave found that in several of these cases journalists were intentionally targeted by the Israeli military -- which constitutes a war crime.
Over the past year IPI has stood with its press freedom partners calling for an immediate end to the killing of journalists in Gaza as well as for international media to be allowed unfettered access to report independently from inside Gaza.
In May, IPI and its partner IMS jointly presented the @ @ @ @ @ @ @ @ @ @ Gaza.
The award recognised the extraordinary courage and resilience that Palestinian journalists have demonstrated in being the world ' s eyes and ears in Gaza.
This week, IPI renewed its call on the international community to protect journalists in Gaza as well as in the West Bank and Lebanon.
Allies of Israel, including Media Freedom Coalition members, must pressure the Israeli government to protect journalist safety and stop attacks on the press.
This also includes the growing media censorship demonstrated by Israel ' s recent closure of Al Jazeera ' s Ramallah bureau.
Raising awareness IPI was at the UN in Geneva this week with its partners Committee to Protect Journalists ( CPJ ), Reporters without Borders ( RSF ), and the European Federation of Journalists ( EFJ ), and others for high-level meetings aimed at raising awareness of the continued attacks on the press and urging the international community to protect journalists.
Among the key messages : The continued killings of journalists in Gaza -- and corresponding impunity -- endangers journalists and press freedom everyone. @ @ @ @ @ @ @ @ @ @ this week ' s Washington Post honours the journalists bravely reporting on the war, often at great personal risk, and underscores IPI ' s solidarity with those that dedicate their lives to uncovering the truth."But it is clear that solidarity is not enough.
Action is needed,"said IPI in its statement."The international community must place effective pressure on the Israeli authorities to comply with international law ; protect the safety of journalists ; investigate the killing of journalists by its forces and secure accountability ; and grant international media outlets immediate and unfettered access to report independently from Gaza."We urge the international community to meet this moment of crisis and stand up for the protection of journalists and freedom of the press in Gaza."An attack against journalists anywhere is an attack against freedom and democracy everywhere."

% matched lemmas: one-year
\end{textsample}
